\documentclass[11pt]{article}
\usepackage{amsmath, amssymb}
\usepackage{hyperref}
\usepackage{graphicx}
\usepackage{listings}
\usepackage{booktabs}


% for code blocks. down below, use \begin{lstlisting} ... \end{lstlisting}
\usepackage{listings}
\usepackage{xcolor}

\lstset{
  language=Java,
  basicstyle=\ttfamily\small,
  keywordstyle=\color{blue},
  commentstyle=\color{gray},
  stringstyle=\color{teal},
  frame=single,
  breaklines=true,
  showstringspaces=false
}

% Define \problemname and \illustration if not using Kattis template
\newcommand{\problemname}[1]{\section*{#1}}
\newcommand{\illustration}[3]{%
    \begin{figure}[h]
        \centering
        \includegraphics[width=#1\textwidth]{#2}
        \caption{#3}
    \end{figure}
}

\begin{document}

\problemname{Preemptive CPU Scheduling}

% Example use of the \illustration command
% \illustration{0.3}{filename.jpg}{Photo by \href{url_here}{link text here, e.g. author name}}

A CPU is a computer processor that does tasks, called processes. The CPU can execute
a single process at a time, but it is also \textit{preemptive}, meaning at any point it can pause
a process to work on another. The CPU is blazingly fast and so switching between 
processes does not take any time. $N$ processes arrive, and the CPU needs to schedule 
the processes. Each process has
\begin{itemize}
    \item An arrival time $a_i$, in units of time
    \item A burst time $b_i$, the total time the process needs to spend on the CPU before it is finished, in units of time
    \item A memory placement time $m_i$, the time taken to place computed results of a finished process in memory, in units of time
\end{itemize}
We consider a process to be \textbf{complete} when it has finished on the CPU and it's result has been placed in memory.
We define the \textbf{waiting time} of a process as the total units of time elapsed between its arrival and its
completion, excluding time the CPU spend working on it. Schedule all $N$ tasks to run on the CPU to find the minimum total
waiting time possible, requiring that all processes must be completed.

\section*{Input}

A single integer $1 \leq N \leq 10^5$, indicating the number of processes, followed by $N$
lines, each containing $0 \leq a_i \leq 10^5$, and $1 \leq b_i, m_i \leq 10^5$

\section*{Output}

A single integer, the minimum possible waiting time to complete all scheduled processes.

\section*{Sample Input \hspace{75px} Sample Output}

\begin{minipage}[t]{0.48\textwidth}
\begin{lstlisting}
3
0 5 2
1 2 1
2 1 3
\end{lstlisting}
\end{minipage}
\hfill
\begin{minipage}[t]{0.48\textwidth}
\begin{lstlisting}
10
\end{lstlisting}
\end{minipage}


\end{document}